\section{Introduction}

An LLM agent does not need to emit malicious text to cause harm. It can send a message to an injected recipient, share the wrong file, turn a draft into a committed action, or carry an attacker-controlled identifier from a tool result into a later call. The security boundary is therefore the point immediately before a structured tool call changes external state. At that boundary, judging only the tool name or the call as one undifferentiated object is too coarse: one call may contain both authorized and unauthorized effects.

We call the missing relation \emph{tool-effect binding}. A tool implementation maps a structured call and pre-state to an external state transition; an effect contract is a conservative security abstraction of that transition. A monitor binds a call when harmless interventions preserve this abstraction, while interventions that change an independently authorizable effect are represented and, when policy requires, change the decision. This semantic layer is upstream of policy enforcement. Capability and least-privilege mechanisms can enforce only what their policy object represents; an omitted Bcc recipient, implicit invitation, or public-link change is invisible to a policy over the enclosing tool name.

The paper develops \sys{}, an architecture with two separated phases. During offline onboarding, an LLM may propose effect templates, but the proposal is only a synthesis hypothesis. Paired sandbox executions intervene on arguments, defaults, and trusted state; the validator compares concrete security-relevant state differences with the atoms produced by the candidate. A concrete effect change with no corresponding atom change is a \emph{mediation gap}; an atom change without a concrete security-effect change is over-sensitivity. Only bounded mappings that survive single-field, interaction, default, expansion, and negative-control tests are frozen. During agent execution, a deterministic pre-commit guard instantiates atoms from the frozen descriptor and totalized call, then compares every atom with a task permission envelope. The runtime does not ask an LLM to reinterpret the contract at each call. We distinguish this normative architecture from our evaluated AgentDojo prototype, which uses one-field state/output differences and therefore does not yet establish complete causal abstraction or an independently grounded envelope.

This structure provides a precise explanation surface: a decision identifies the contract template, concrete atom, value source, matching envelope clause, and mismatch reason. It also exposes the assumptions needed for a safety argument. If contracts conservatively cover true effects, the envelope does not expand from untrusted observations, every effectful call is mediated, checked values are the values executed, and uncertainty fails closed, then every committed effect remains within the original task envelope. The same argument extends by induction to every prefix of a multi-step trajectory.

Our evaluation protocol follows the obligations rather than relying on one aggregate benchmark score. Counterfactual registration tests contract coverage and field relevance. Same-protocol AgentDojo runs test utility and indirect-prompt-injection attack success under one victim checkpoint. Stateful long-horizon tasks, controlled ablations, and adaptive attacks test legitimate resolution, prefix behavior, representation granularity, provenance, envelope planning, replan handling, and fail-closed behavior. We report only completed rows whose task manifests and row-level artifacts pass the common-protocol gates.

The intended contributions are:
\begin{itemize}[leftmargin=*]
  \item a pre-commit formulation of tool-effect binding and a concrete effect-atom mediation object;
  \item counterfactual contract registration that tests whether candidate atoms mediate security-relevant execution differences, with explicit invariance, mediation-gap, and interaction witnesses;
  \item a deterministic runtime that binds frozen descriptors to task-scoped authority and emits auditable decision witnesses;
  \item a conditional single-call and long-trajectory confinement argument with explicit trusted-computing-base obligations; and
  \item an obligation-driven evaluation methodology, with completed official AgentDojo and implementation-level mediation evidence and common-protocol long-horizon, adaptive, ablation, and overhead gates.
\end{itemize}
